\documentclass[aspectratio=169]{beamer}

\usepackage[T1]{fontenc}
\usepackage[brazil]{babel}
\usepackage[utf8]{inputenc}
\usepackage{lmodern}
\usepackage{booktabs}
\usepackage{array}
\usepackage{amsmath}
\usepackage{tikz}
\usetikzlibrary{arrows.meta, positioning}

\IfFileExists{beamerthemeInf.sty}{
  \usetheme{Inf}
}{
  \usetheme{Madrid}
}

\providecommand{\InfTitlePage}{
  \begin{frame}[plain]
    \titlepage
  \end{frame}
}

\title{TSDAE ainda é útil para recuperação semântica em português?}
\subtitle{Uma análise experimental com BERT e embeddings multilíngues modernos}



\institute{
Instituto de Informática --- UFRGS
}

\newcommand{\model}[1]{\texttt{#1}}

\begin{document}

\InfTitlePage

\begin{frame}
  \frametitle{Agenda}
  \tableofcontents
\end{frame}

% ==========================================================
\section{Motivação}

\begin{frame}
\frametitle{Contexto}

\begin{block}{Recuperação semântica}
Em tarefas de recuperação de informação, o objetivo é ordenar documentos de acordo com sua relevância para uma consulta.

\vspace{0.15cm}

Modelos densos de embeddings permitem recuperar documentos semanticamente relacionados, mesmo quando não há correspondência lexical direta entre consulta e documento.
\end{block}

\vspace{0.25cm}

\begin{block}{Mudança recente}
Modelos multilíngues modernos, como E5 e BGE-M3, já apresentam forte desempenho em diferentes línguas e domínios.

\vspace{0.15cm}

Isso levanta uma dúvida: a adaptação não supervisionada com TSDAE ainda traz ganho real quando aplicada a modelos BERT?
\end{block}

\end{frame}

\begin{frame}
\frametitle{TSDAE: Relembrando}

\begin{block}{Ideia central}
O \textit{Transformer-based Sequential Denoising Auto-Encoder} treina um modelo para reconstruir uma sentença original a partir de uma versão corrompida.

\vspace{0.2cm}

Durante esse processo, o encoder é incentivado a produzir representações que preservam informação semântica relevante.
\end{block}

\vspace{0.25cm}

\[
x
\;\xrightarrow{\text{ruído}}\;
\tilde{x}
\;\xrightarrow{\text{encoder}}\;
z
\;\xrightarrow{\text{decoder}}\;
\hat{x}
\]

\vspace{0.15cm}

\begin{alertblock}{Uso no experimento}
Após a adaptação, o encoder BERT é usado para gerar embeddings de consultas e documentos.
\end{alertblock}

\end{frame}

\begin{frame}
\frametitle{Pergunta de Pesquisa}

\begin{alertblock}{Questão central}
A adaptação não supervisionada com TSDAE em modelos BERT ainda gera embeddings úteis para recuperação semântica em português quando comparada a embeddings multilíngues modernos já prontos?
\end{alertblock}

\vspace{0.3cm}

\textbf{Hipóteses investigadas}

\begin{itemize}
  \item BERT + TSDAE pode superar o BERT original;
  \item O ganho depende da estratégia de ruído usada no TSDAE;
  \item Modelos modernos como E5 e BGE-M3 podem superar o BERT adaptado isoladamente;
  \item Mesmo assim, BERT + TSDAE pode recuperar documentos relevantes diferentes;
  \item Essa complementaridade pode ser útil em fusões de rankings.
\end{itemize}

\end{frame}

% ==========================================================
\section{Desenho Experimental}

\begin{frame}
\frametitle{Visão Geral do Experimento}

\begin{center}
\begin{tikzpicture}[
  node distance=0.45cm,
  box/.style={
    draw,
    rounded corners,
    align=center,
    minimum width=4.8cm,
    minimum height=0.65cm
  },
  arrow/.style={-{Latex[length=2mm]}, thick}
]

\node[box] (data) {Dataset Quati\\ corpus, queries e qrels};
\node[box, below=of data] (systems) {Sistemas avaliados\\ BM25, embeddings prontos, BERT e BERT + TSDAE};
\node[box, below=of systems] (retrieval) {Geração de rankings\\ busca lexical e busca vetorial};
\node[box, below=of retrieval] (eval) {Avaliação quantitativa\\ MRR, Recall@k, Hit@k e nDCG@k};
\node[box, below=of eval] (analysis) {Análises adicionais\\ complementaridade, fusão e bootstrap};

\draw[arrow] (data) -- (systems);
\draw[arrow] (systems) -- (retrieval);
\draw[arrow] (retrieval) -- (eval);
\draw[arrow] (eval) -- (analysis);

\end{tikzpicture}
\end{center}

\vspace{0.15cm}

\begin{block}{Foco metodológico}
Todos os métodos são comparados sobre o mesmo conjunto de consultas, documentos e julgamentos de relevância.
\end{block}

\end{frame}

\begin{frame}
\frametitle{Unidade de Comparação}

\textbf{Cada sistema gera um ranking para cada consulta.}

\vspace{0.3cm}

\begin{columns}

\column{0.48\textwidth}
\begin{block}{Entrada}
\begin{itemize}
  \item Consulta em português;
  \item Corpus de documentos/passagens;
  \item Modelo ou método de recuperação.
\end{itemize}
\end{block}

\column{0.48\textwidth}
\begin{block}{Saída}
\begin{itemize}
  \item Lista ordenada de documentos;
  \item Top-$k$ documentos recuperados;
  \item Scores ou posições no ranking.
\end{itemize}
\end{block}

\end{columns}

\vspace{0.3cm}

\begin{alertblock}{Comparação}
A avaliação é feita por consulta e depois agregada para obter métricas globais.
\end{alertblock}

\end{frame}

% ==========================================================
\section{Dataset}

\begin{frame}
\frametitle{Dataset Utilizado: Quati}

\begin{block}{Motivação para a escolha}
O Quati é utilizado por ser um dataset de recuperação de informação em português brasileiro, com consultas, corpus e julgamentos de relevância.
\end{block}

\vspace{0.25cm}

\begin{table}
\centering
\small
\begin{tabular}{@{}ll@{}}
\toprule
\textbf{Componente} & \textbf{Uso no experimento} \\
\midrule
Queries & Consultas usadas para gerar os rankings \\
Corpus & Documentos/passagens candidatos à recuperação \\
Qrels & Julgamentos de relevância usados apenas na avaliação \\
\bottomrule
\end{tabular}
\end{table}

\vspace{0.25cm}

\begin{alertblock}{Ponto importante}
Os qrels não são usados para adaptar o BERT com TSDAE; eles entram somente na etapa de avaliação.
\end{alertblock}

\end{frame}

\begin{frame}
\frametitle{Preparação dos Dados}

\textbf{Organização do experimento}

\begin{itemize}
  \item Padronização dos identificadores de queries e documentos;
  \item Separação entre corpus, consultas e qrels;
  \item Uso do corpus textual para indexação e adaptação não supervisionada;
  \item Uso das queries para gerar rankings;
  \item Uso dos qrels para calcular métricas de recuperação.
\end{itemize}

\vspace{0.3cm}

\begin{block}{Controle contra vazamento de informação}
A adaptação TSDAE utiliza textos sem rótulos de relevância. As informações de relevância são reservadas para a avaliação final.
\end{block}

\end{frame}

% ==========================================================
\section{Sistemas Avaliados}

\begin{frame}
\frametitle{Grupos de Sistemas Avaliados}

\begin{table}
\centering
\small
\begin{tabular}{@{}p{0.28\textwidth}p{0.62\textwidth}@{}}
\toprule
\textbf{Grupo} & \textbf{Papel no experimento} \\
\midrule
BM25 & Baseline lexical baseado em correspondência de termos \\
Embeddings multilíngues prontos & Baselines densos fortes, sem adaptação adicional \\
BERT original & Modelo base usado como referência antes da adaptação \\
BERT + TSDAE & Modelo BERT adaptado de forma não supervisionada com diferentes ruídos \\
Fusão de rankings & Combinação de rankings para avaliar complementaridade \\
\bottomrule
\end{tabular}
\end{table}

\vspace{0.25cm}

\begin{alertblock}{Ponto central}
O TSDAE não é aplicado a todos os baselines densos. Ele é aplicado ao modelo BERT, e os modelos modernos servem como referência competitiva.
\end{alertblock}

\end{frame}

\begin{frame}
\frametitle{Baselines Densos Multilíngues}

\textbf{Modelos usados sem adaptação adicional}

\vspace{0.2cm}

\begin{itemize}
  \item multilingual-E5;
  \item BGE-M3;
  \item LaBSE;
  \item MiniLM Multilingual;
  \item DistilUSE.
\end{itemize}

\vspace{0.35cm}

\begin{block}{Função desses modelos}
Eles representam o estado atual de embeddings multilíngues prontos para uso, servindo como referência para avaliar se o BERT adaptado com TSDAE ainda é competitivo.
\end{block}

\vspace{0.15cm}



\end{frame}

\begin{frame}
\frametitle{Modelo Adaptado com TSDAE}

\textbf{Modelo base: BERT}

\vspace{0.25cm}

\begin{itemize}
  \item O BERT é usado como modelo base para a adaptação não supervisionada;
  \item Primeiro, avalia-se o BERT original, sem TSDAE;
  \item Em seguida, avaliam-se versões do BERT adaptadas com TSDAE;
  \item Cada versão adaptada usa uma estratégia diferente de corrupção textual.
\end{itemize}

\vspace{0.3cm}

\begin{block}{Comparação principal}
BERT original vs. BERT + TSDAE.

\vspace{0.15cm}

Depois, compara-se o comportamento dessas versões com modelos modernos como E5 e BGE-M3.
\end{block}

\end{frame}

\begin{frame}
\frametitle{Baselines Lexicais}

\textbf{BM25}

\vspace{0.2cm}

\begin{itemize}
  \item Recuperação baseada em termos;
  \item Forte baseline clássico para busca textual;
  \item Não depende de embeddings;
  \item Serve como comparação contra os métodos densos.
\end{itemize}

\vspace{0.3cm}

\begin{block}{Por que incluir BM25?}
Um modelo denso só é realmente útil se superar, complementar ou combinar bem com um baseline lexical forte.
\end{block}

\end{frame}

% ==========================================================
\section{Adaptação com TSDAE}

\begin{frame}
\frametitle{TSDAE no Experimento}

\textbf{Etapa de adaptação não supervisionada}

\vspace{0.25cm}

\begin{enumerate}
  \item Seleciona-se o modelo BERT como base;
  \item Sentenças ou passagens do corpus são corrompidas;
  \item O modelo é treinado para reconstruir o texto original;
  \item O encoder BERT adaptado é usado para gerar embeddings;
  \item Os embeddings são avaliados em recuperação semântica.
\end{enumerate}

\vspace{0.25cm}

\begin{alertblock}{Ideia experimental}
A adaptação é tratada como uma intervenção: comparamos o BERT antes e depois do TSDAE.
\end{alertblock}

\end{frame}

\begin{frame}
\frametitle{Configurações de TSDAE}

O TSDAE é aplicado ao modelo BERT usando diferentes estratégias de ruído.

\vspace{0.25cm}

\begin{table}
\centering
\small
\begin{tabular}{@{}ll@{}}
\toprule
\textbf{Configuração} & \textbf{Descrição} \\
\midrule
\texttt{delete30} & Remove 30\% dos tokens \\
\texttt{delete60} & Remove 60\% dos tokens \\
\texttt{mask30} & Mascara 30\% dos tokens \\
\texttt{span\_delete30} & Remove spans contínuos do texto \\
\bottomrule
\end{tabular}
\end{table}

\vspace{0.3cm}

\begin{alertblock}{Pergunta metodológica}
Qual tipo de ruído produz embeddings BERT mais úteis para recuperação semântica em português?
\end{alertblock}

\end{frame}

\begin{frame}
\frametitle{Por que Avaliar Diferentes Ruídos?}

\begin{columns}

\column{0.48\textwidth}
\begin{block}{Ruídos mais leves}
\begin{itemize}
  \item Preservam mais informação lexical;
  \item Podem facilitar a reconstrução;
  \item Podem gerar mudanças menores no modelo.
\end{itemize}
\end{block}

\column{0.48\textwidth}
\begin{block}{Ruídos mais fortes}
\begin{itemize}
  \item Forçam maior abstração semântica;
  \item Podem melhorar robustez;
  \item Também podem degradar o sinal original.
\end{itemize}
\end{block}

\end{columns}

\vspace{0.3cm}

\begin{alertblock}{Questão metodológica}
A avaliação não considera apenas se TSDAE melhora o BERT, mas qual configuração altera o ranking de forma útil.
\end{alertblock}

\end{frame}

% ==========================================================
\section{Pipeline de Recuperação}

\begin{frame}
\frametitle{Pipeline de Recuperação Densa}

\begin{columns}

\column{0.48\textwidth}
\begin{block}{Indexação}
\begin{itemize}
  \item O corpus é codificado pelo modelo;
  \item Cada documento recebe um embedding;
  \item Os vetores são armazenados em um índice FAISS.
\end{itemize}
\end{block}

\column{0.48\textwidth}
\begin{block}{Consulta}
\begin{itemize}
  \item A query é codificada pelo mesmo modelo;
  \item Busca-se os vetores mais próximos;
  \item O sistema retorna os top-$k$ documentos.
\end{itemize}
\end{block}

\end{columns}

\vspace{0.3cm}

\begin{alertblock}{Controle experimental}
O mesmo procedimento de recuperação é usado para modelos densos prontos, BERT original e BERT adaptado.
\end{alertblock}

\end{frame}

\begin{frame}
\frametitle{Rankings Gerados}

Para cada consulta, cada sistema produz uma lista ordenada de documentos.

\vspace{0.25cm}

\begin{table}
\centering
\small
\begin{tabular}{@{}lll@{}}
\toprule
\textbf{Sistema} & \textbf{Representação} & \textbf{Saída} \\
\midrule
BM25 & Termos e frequências & Ranking lexical \\
Embeddings multilíngues prontos & Embeddings congelados & Ranking semântico \\
BERT original & Embeddings BERT sem adaptação & Ranking semântico \\
BERT + TSDAE & Embeddings BERT adaptados & Ranking semântico adaptado \\
Fusão & Combinação de rankings & Ranking combinado \\
\bottomrule
\end{tabular}
\end{table}

\vspace{0.25cm}

\begin{block}{Base da análise}
As métricas e análises de complementaridade são calculadas a partir desses rankings.
\end{block}

\end{frame}

% ==========================================================
\section{Métricas de Avaliação}

\begin{frame}
\frametitle{Métricas Tradicionais}

\begin{table}
\centering
\small
\begin{tabular}{@{}ll@{}}
\toprule
\textbf{Métrica} & \textbf{O que mede} \\
\midrule
MRR & Posição do primeiro documento relevante recuperado \\
Recall@k & Fração dos documentos relevantes presentes no top-$k$ \\
Hit@k & Se há ao menos um documento relevante no top-$k$ \\
nDCG@k & Qualidade do ranking considerando posição e relevância \\
\bottomrule
\end{tabular}
\end{table}

\vspace{0.3cm}

\begin{alertblock}{Interpretação}
As métricas capturam aspectos diferentes: encontrar rápido o primeiro relevante, cobrir relevantes e ordenar bem o ranking.
\end{alertblock}

\end{frame}

\begin{frame}
\frametitle{Avaliação por Consulta}

Além das médias globais, o experimento analisa o comportamento por query.

\vspace{0.25cm}

\begin{itemize}
  \item Quais consultas melhoram após TSDAE;
  \item Quais consultas pioram após TSDAE;
  \item Em quais casos modelos modernos já são suficientes;
  \item Em quais casos BM25 recupera documentos que os modelos densos não recuperam;
  \item Em quais casos BERT + TSDAE recupera documentos relevantes exclusivos.
\end{itemize}

\vspace{0.3cm}

\begin{block}{Motivação}
Duas abordagens podem ter médias parecidas, mas errar e acertar em consultas diferentes.
\end{block}

\end{frame}

\begin{frame}
\frametitle{Significância Estatística}

\textbf{Comparações pareadas entre sistemas}

\vspace{0.25cm}

\begin{itemize}
  \item A unidade de reamostragem é a query;
  \item Calcula-se a diferença de desempenho entre dois sistemas;
  \item Bootstrap é usado para estimar intervalos de confiança;
  \item O foco é avaliar a incerteza das diferenças observadas.
\end{itemize}

\vspace{0.3cm}

\begin{alertblock}{Exemplo de comparação}
BERT original vs. BERT + TSDAE.

BM25 vs. modelo denso.

Modelo multilíngue moderno vs. BERT + TSDAE.

Melhor modelo individual vs. fusão de rankings.
\end{alertblock}

\end{frame}

% ==========================================================
\section{Complementaridade e Fusão}

\begin{frame}
\frametitle{Complementaridade entre Rankings}

\begin{block}{Ideia}
Um modelo pode ser útil mesmo sem ter a melhor métrica média, desde que recupere documentos relevantes diferentes dos demais modelos.
\end{block}

\vspace{0.25cm}

\textbf{Análises possíveis}

\begin{itemize}
  \item Sobreposição entre top-$k$ documentos;
  \item Documentos relevantes exclusivos por modelo;
  \item Consultas em que um modelo acerta e outro erra;
  \item Diferença entre rankings gerados por BM25, embeddings prontos, BERT original e BERT + TSDAE.
\end{itemize}

\vspace{0.25cm}

\begin{alertblock}{Pergunta central}
BERT + TSDAE gera rankings redundantes ou complementares aos modelos modernos?
\end{alertblock}

\end{frame}

\begin{frame}
\frametitle{Fusão de Rankings}

Duas estratégias de ensemble são consideradas.

\vspace{0.25cm}

\begin{columns}

\column{0.48\textwidth}
\begin{block}{Score Fusion}
\begin{itemize}
  \item Normaliza scores;
  \item Combina pontuações;
  \item Gera um ranking médio.
\end{itemize}
\end{block}

\column{0.48\textwidth}
\begin{block}{Reciprocal Rank Fusion}
\begin{itemize}
  \item Usa posições no ranking;
  \item Não exige scores comparáveis;
  \item Favorece documentos bem ranqueados por múltiplos sistemas.
\end{itemize}
\end{block}

\end{columns}

\vspace{0.25cm}

\[
\text{RRF}(d) = \sum_{m \in M} \frac{1}{c + \text{rank}_m(d)}
\]

\end{frame}

\begin{frame}
\frametitle{Por que Ensembles?}

\begin{block}{Motivação}
Se BM25, modelos multilíngues modernos e BERT adaptado com TSDAE recuperam documentos relevantes diferentes, a combinação pode superar os modelos isolados.
\end{block}

\vspace{0.25cm}

\textbf{O experimento avalia se:}

\begin{itemize}
  \item BERT + TSDAE contribui para ensembles;
  \item A fusão melhora métricas tradicionais;
  \item A melhora vem de complementaridade real ou apenas de reforço dos mesmos documentos;
  \item O BERT adaptado pode agregar valor quando combinado com modelos fortes.
\end{itemize}

\end{frame}

% ==========================================================
\section{Cuidados Experimentais}

\begin{frame}
\frametitle{Controles Metodológicos}

\begin{itemize}
  \item Mesmas queries para todos os sistemas;
  \item Mesmo corpus de documentos candidatos;
  \item Mesmos qrels para avaliação;
  \item Comparação direta entre BERT original e BERT adaptado;
  \item Avaliação dos rankings no mesmo top-$k$;
  \item Separação entre adaptação não supervisionada e avaliação supervisionada.
\end{itemize}

\vspace{0.3cm}

\begin{alertblock}{Objetivo}
Garantir que diferenças observadas sejam atribuídas ao método de recuperação ou à adaptação, não a mudanças no protocolo experimental.
\end{alertblock}

\end{frame}

\begin{frame}
\frametitle{O que a Metodologia Permite Responder}

Ao final da avaliação, o desenho experimental permite investigar:

\vspace{0.2cm}

\begin{enumerate}
  \item TSDAE melhora o BERT em recuperação semântica em português?
  \item O ganho depende do tipo de ruído usado?
  \item Modelos modernos como E5 e BGE-M3 reduzem a necessidade de TSDAE?
  \item BERT + TSDAE ainda recupera documentos relevantes complementares?
  \item BERT + TSDAE ainda é útil em ensembles ou cenários de complementaridade?
\end{enumerate}

\vspace{0.3cm}

\begin{block}{Foco do estudo}
Não apenas medir se TSDAE ganha, mas entender quando ele ainda vale a pena.
\end{block}

\end{frame}

\begin{frame}
\frametitle{Resumo da Metodologia}

\begin{center}
\begin{tikzpicture}[
  node distance=0.35cm,
  box/.style={
    draw,
    rounded corners,
    align=center,
    minimum width=5.3cm,
    minimum height=0.58cm
  },
  arrow/.style={-{Latex[length=2mm]}, thick}
]

\node[box] (q) {Quati: queries, corpus e qrels};
\node[box, below=of q] (b) {BM25, embeddings multilíngues prontos e BERT original};
\node[box, below=of b] (t) {Adaptação do BERT com TSDAE e diferentes ruídos};
\node[box, below=of t] (r) {Geração de rankings via BM25 e FAISS};
\node[box, below=of r] (m) {MRR, Recall@k, Hit@k e nDCG@k};
\node[box, below=of m] (c) {Complementaridade, fusão e bootstrap};

\draw[arrow] (q) -- (b);
\draw[arrow] (b) -- (t);
\draw[arrow] (t) -- (r);
\draw[arrow] (r) -- (m);
\draw[arrow] (m) -- (c);

\end{tikzpicture}
\end{center}

\end{frame}

\begin{frame}
\frametitle{Próxima Parte da Apresentação}

\begin{block}

\begin{itemize}
  \item desempenho dos baselines;
  \item impacto da adaptação TSDAE no BERT;
  \item comparação entre ruídos;
  \item análise por query;
  \item complementaridade entre rankings;
  \item desempenho dos ensembles;
  \item significância estatística.
\end{itemize}
\end{block}

\vspace{0.25cm}

\begin{alertblock}{Obrigado!}

\end{alertblock}

\end{frame}

\end{document}